\section{ARCore платформа}

\textit{ARCore} је платформа компаније \textit{Google} за развој апликација проширене стварности на \textit{Android} уређајима. Реч је о бесплатном комплету за развој софтвера (\textit{SDK}) који програмерима ставља на располагање готова решења за најтеже проблеме проширене стварности: одређивање положаја уређаја у простору, препознавање површина у окружењу и процену осветљења сцене. Платформа се на уређаје испоручује кроз системску компоненту \textit{Google Play Services for AR}, која се аутоматски ажурира путем продавнице \textit{Google Play}, па апликације увек користе актуелну верзију без потребе да је саме укључују у инсталациони пакет. \textit{ARCore} је подржан на стотинама сертификованих модела телефона и таблета различитих произвођача, при чему сертификација гарантује да су камера, сензори покрета и процесор уређаја довољно квалитетни за стабилно искуство проширене стварности. Захваљујући томе, једна иста апликација ради на огромном броју уређаја који се већ налазе у рукама корисника, без икаквог додатног хардвера.

Суштина платформе може се сажети у три темељне способности: праћење покрета, разумевање окружења и процена осветљења. Свака од њих решава по један предуслов уверљиве проширене стварности, а заједно омогућавају да виртуелни садржај делује као да заиста постоји у физичком простору.

\subsection{Праћење покрета}

Да би виртуелни објекат остао на свом месту док се корисник креће, телефон у сваком тренутку мора да зна где се налази и куда је усмерен. \textit{ARCore} тај проблем решава поступком који се назива визуелно-инерцијална одометрија. Платформа на слици са камере проналази визуелно упечатљиве тачке --- ивице, углове и контрастне детаље текстура --- и прати како се њихов положај мења из кадра у кадар. Истовремено очитава податке са акцелерометра и жироскопа, сензора који мере убрзање и ротацију уређаја. Комбиновањем ова два извора информација, који се међусобно допуњују и исправљају, \textit{ARCore} израчунава положај и оријентацију телефона у простору са шест степени слободе: три транслације дуж просторних оса и три ротације око њих. Захваљујући томе, виртуелни садржај се у сваком кадру исцртава из тачно одговарајуће перспективе --- када корисник приђе виртуелном објекту, објекат се увећава; када га заобиђе, види га са друге стране, баш као да је реч о стварном предмету.

\subsection{Разумевање окружења}

Поред сопственог положаја, уређај мора да разуме и простор око себе. \textit{ARCore} непрекидно анализира распоред карактеристичних тачака које прати и, када уочи групе тачака које леже у истој равни, закључује да је пронашао површину: под, радни сто, седиште столице или зид. Платформа препознаје и хоризонталне и вертикалне равни, одређује њихове границе и проширује их како камера обухвата нове делове простора. Тако изграђен модел окружења омогућава да се виртуелни објекти постављају на смислена места --- фигура на под, а слика на зид --- уместо да лебде у ваздуху.

Важан механизам који овај модел чини употребљивим из перспективе корисника јесте тестирање погодака (енгл. \textit{hit-testing}). Када корисник додирне екран, \textit{ARCore} из положаја камере кроз додирнуту тачку на екрану пушта замишљени зрак у тродимензионални простор и израчунава где тај зрак сече препознате равни или друге елементе окружења. Резултат је прецизна тачка у простору, са положајем и оријентацијом, на коју апликација може да постави виртуелни садржај. Овај механизам је темељ великог броја АР апликација, јер сваки додир корисника претвара у тачку стварног простора на коју се може везати дигитални садржај.

\subsection{Процена осветљења}

Трећа способност платформе односи се на светлост. Виртуелни објекат исцртан без обзира на услове осветљења у просторији одмах одаје своју вештачку природу: превише је светао у мрачној соби или безличан на јаком сунцу. \textit{ARCore} зато из слике са камере процењује карактеристике осветљења сцене --- просечан интензитет, боју светлости, а у режиму \textit{Environmental HDR} и правац главног извора света те околинско осветљење погодно за реалистичне одсјаје. Апликација те податке користи приликом сенчења виртуелних објеката, па они добијају сенке и одсјаје усклађене са стварним окружењем, што значајно доприноси утиску да објекат заиста припада сцени.

\subsection{Дубина, сидра и објекти праћења}

Поред три основне способности, за практичан рад важно је још неколико појмова. \textit{Depth API} омогућава да \textit{ARCore}, користећи само једну обичну камеру, за сваки пиксел слике процени удаљеност од уређаја и тако изгради дубинску мапу сцене. Захваљујући њој, виртуелни садржај се може постављати на произвољне површине, а не само на препознате равни, док виртуелни објекти могу бити заклоњени стварним предметима који се налазе испред њих (оклузија), што додатно појачава уверљивост приказа.

Виртуелни садржај се за стварни свет везује помоћу сидара (енгл. \textit{anchor}). Сидро представља фиксну позу у простору коју платформа активно одржава: како \textit{ARCore} временом прикупља нове информације и прецизира своје разумевање сцене, он коригује и положаје сидара, па садржај причвршћен за њих остаје „прикован" за исто место у стварном свету уместо да постепено отплута. Сидра се по правилу везују за објекте праћења (енгл. \textit{trackable}) --- равни, карактеристичне тачке и друге елементе које платформа препознаје и прати кроз време.

У пракси се \textit{ARCore} користи или непосредно, кроз његов изворни програмски интерфејс, или кроз библиотеке вишег нивоа које поједностављују исцртавање тродимензионалног садржаја. Међу њима се издваја \textit{SceneView}, библиотека која обједињује \textit{ARCore} и графички погон \textit{Filament}, па програмер уместо ниског графичког кода ради са сценом, чворовима и моделима, што знатно убрзава развој.

\subsection{Праћење слика (Augmented Images)}

Поред равни и карактеристичних тачака, \textit{ARCore} уме да препознаје и прати унапред задате дводимензионалне слике у стварном свету --- способност позната под називом \textit{Augmented Images} (праћење слика). Програмер унапред саставља базу слика, објекат типа \texttt{AugmentedImageDatabase}, у коју сваку референтну слику додаје позивом \texttt{addImage} уз јединствено име и, опционо, физичку ширину одштампаног примерка изражену у метрима. Када је база придружена конфигурацији сесије, \textit{ARCore} у сваком кадру упоређује карактеристичне тачке слике са камере са тачкама слика из базе. Чим препозна неку од њих, платформа враћа објекат \texttt{AugmentedImage} који носи име препознате слике, њену позу у простору (\texttt{centerPose}), процењене димензије (\texttt{extentX} и \texttt{extentZ}) и стање праћења. Препозната слика тиме постаје пуноправан објекат праћења: за њу се може везати сидро, а виртуелни садржај постављен у њеном координатном систему прати је док се корисник креће.

Квалитет препознавања пресудно зависи од саме слике маркера. Добра референтна слика богата је детаљима, има висок контраст и избегава понављајуће шаре и велике једнобојне површине, јер се препознавање заснива на распореду карактеристичних тачака. \textit{Google} уз платформу испоручује алат \texttt{arcoreimg}, који слику оцењује на скали од 0 до 100; препоручује се оцена од најмање 75 да би праћење било поуздано. Навођење физичке ширине слике додатно је важно: када зна стварне димензије маркера, платформа одмах по препознавању тачно процењује његову удаљеност и позу, уместо да их постепено прецизира посматрањем из више углова. Једна база може да садржи и до хиљаду слика, при чему се истовремено прати до двадесет препознатих маркера.

За апликацију из овог рада праћење слика представља природан избор и средишњи механизам читавог решења. Одштампани маркер на вратима учионице једнозначно идентификује просторију --- име препознате слике у себи носи ознаку учионице, па апликација одмах зна од ког сервера затражити податке --- док сидро везано за позу слике обезбеђује да информациони панел остане „прикован" изнад врата и када се корисник помера, приближава или мења угао гледања. Поврх свега, платформа ради на широком спектру постојећих уређаја без додатне опреме, чиме апликација постаје доступна најширем кругу корисника. Начин на који се ове могућности конкретно користе у развоју приказан је у наредним поглављима.
